\section{AHSME 1965}
        \begin{enumerate}
        	\item (\textit{AHSME 1965}) The number of real values of $x$ satisfying the equation $2^{2x^2 - 7x + 5} = 1$ is: 


 $\textbf{(A)}\ 0 \qquad  \textbf{(B) }\ 1 \qquad  \textbf{(C) }\ 2 \qquad  \textbf{(D) }\ 3 \qquad  \textbf{(E) }\ \text{more than 4}$


 
 \textbf{Answer: }C



\textbf{Solution 1}

Solution by e\_power\_pi\_times\_i

 Take the logarithm with a base of $2$ to both sides, resulting in the equation $2x^2-7x+5 = 0$. Factoring results in $(2x-5)(x-1) = 0$, so there are $\boxed{\textbf{(C) } 2}$ real solutions.

 



\textbf{Solution 2}

Notice that $a^0=1, a>0$. So $2^0=1$. So $2x^2-7x+5=0$. Evaluating the \href{https://artofproblemsolving.com/wiki/index.php?title=Discriminant}{discriminant}, we see that it is equal to $7^2-4*2*5=9$. So this means that the equation has two real solutions. Therefore, select $\fbox{C}$.

 


	\item (\textit{AHSME 1965}) A regular hexagon is inscribed in a circle. The ratio of the length of a side of the hexagon to the length of the 
shorter of the arcs intercepted by the side, is: 


 $\textbf{(A)}\ 1: 1 \qquad  \textbf{(B) }\ 1: 6 \qquad  \textbf{(C) }\ 1: \pi \qquad  \textbf{(D) }\ 3: \pi \qquad  \textbf{(E) }\ 6:\pi$


 
 \textbf{Answer: }D



\textbf{Solution}

Suppose that each side of the hexagon is $3$. Then the distance from each vertex of the hexagon to the center is also $3$, so that the circle has radius $3$. Since the circle has circumference $2\pi(3) = 6\pi$, the arc intercepted by any side (which measures $60^\circ$) has length $\frac{1}{6}*6\pi = \pi$, so that our answer is $\boxed{\text{(D)}}$ and we are done.

 


	\item (\textit{AHSME 1965}) The expression $(81)^{ - 2^{ - 2}}$ has the same value as: 


 $\textbf{(A)}\ \frac {1}{81} \qquad  \textbf{(B) }\ \frac {1}{3} \qquad  \textbf{(C) }\ 3 \qquad  \textbf{(D) }\ 81\qquad \textbf{(E) }\ 81^4$


 
 \textbf{Answer: }B



\textbf{Solution}

Let us recall $\text{PEMDAS}$. We  calculate the exponent first. $(-2)^{-2}=\frac{1}{(-2)^2}=\frac{1}{4}$ When we substitute, we get $81^{\frac{1}{4}}=\sqrt[4]{81}=\boxed{\textbf{(C) }3}$.

 


	\item (\textit{AHSME 1965}) Line $\ell_2$ intersects line $\ell_1$ and line $\ell_3$ is parallel to $\ell_1$. The three lines are distinct and lie in a plane. 
The number of points equidistant from all three lines is: 


 $\textbf{(A)}\ 0 \qquad  \textbf{(B) }\ 1 \qquad  \textbf{(C) }\ 2 \qquad  \textbf{(D) }\ 4 \qquad  \textbf{(E) }\ 8$


 
 \textbf{Answer: }C



\textbf{Solution}


                    \begin{figure}[H]
                        centering
                        \includegraphics[width=0.5\textwidth]{Solution/AHSME_I_1965_4_Solution_Figure_1.png}
                    \end{figure}
                    

 The lines are coplanar, $\ell_1 \parallel \ell_3$, and $\ell_1$ intersects $\ell_2$. Therefore, $\ell_3$ also intersects $\ell_2$. The  \href{https://artofproblemsolving.com/wiki/index.php?title=Locus}{locus} of all points equidistant from parallel lines $\ell_1$ and $\ell_3$ is a third parallel line in between them. Let this line be $\ell_4$, and let the distance from $\ell_4$ to either $\ell_1$ or $\ell_3$ be $d$. The points equidistant from lines $\ell_1$, $\ell_2$, and $\ell_3$ must all lie on $\ell_4$ and be a distance $d$ from line $\ell_2$. There are only 2 points, on either side of $\ell_2$, which satisfy these conditions. Thus, our answer is $\fbox{(C) 2}$.

 


	\item (\textit{AHSME 1965}) When the repeating decimal $0.363636\ldots$ is written in simplest fractional form, the sum of the numerator and denominator is: 


 $\textbf{(A)}\ 15 \qquad  \textbf{(B) }\ 45 \qquad  \textbf{(C) }\ 114 \qquad  \textbf{(D) }\ 135 \qquad  \textbf{(E) }\ 150$


 
 \textbf{Answer: }A



\textbf{Solution}

We let $x=0.\overline{36}$. Thus, $100x=36.\overline{36}$. We find that $100x-x=99x=36.\overline{36}-0.\overline{36}=36$, or $x=\frac{36}{99}=\frac{4}{11}$. Since $4+11=15$, the answer is $\boxed{\textbf{(A)}\ 15}$.

 


	\item (\textit{AHSME 1965}) If $10^{\log_{10}9} = 8x + 5$ then $x$ equals: 


 $\textbf{(A)}\ 0 \qquad  \textbf{(B) }\ \frac {1}{2} \qquad  \textbf{(C) }\ \frac {5}{8} \qquad  \textbf{(D) }\ \frac{9}{8}\qquad \textbf{(E) }\ \frac{2\log_{10}3-5}{8}$


 
 \textbf{Answer: }B



\textbf{Solution}

Notice that $10^{\log_{10} 9} = 9$. Therefore, the condition we are looking for is $9=8x+5$, or $x=\frac{1}{2}$. $\text{So the answer is }  \boxed{\textbf{(B)}}$

 


	\item (\textit{AHSME 1965}) The sum of the reciprocals of the roots of the equation $ax^2 + bx + c = 0$ is: 


 $\textbf{(A)}\ \frac {1}{a} + \frac {1}{b} \qquad  \textbf{(B) }\ - \frac {c}{b} \qquad  \textbf{(C) }\ \frac{b}{c}\qquad \textbf{(D) }\ -\frac{a}{b}\qquad \textbf{(E) }\ -\frac{b}{c}$


 
 \textbf{Answer: }E



\textbf{Solution 1}

Using \href{https://artofproblemsolving.com/wiki/index.php?title=Vieta\%27s\_formulas}{Vieta's formulas}, we can write the sum of the roots of any quadratic equation in the form $ax^2+bx+c = 0$ as $\frac{-b}{a}$, and the product as $\frac{c}{a}$.

 If $r$ and $s$ are the roots, then the sum of the reciprocals of the roots is $\frac{1}{r} + \frac{1}{s} = \frac{r+s}{rs}$.

 Applying the formulas, we get $\frac{\frac{-b}{a}}{\frac{c}{a}}$, or $\frac {-b}{c}$ => $\boxed{E}$.

 



\textbf{Solution 2}

Let $f(x)=ax^2+bx+c$. $f(x)$ has roots $r$ and $s$, and so $x^2f(\frac{1}{x})$ has roots $\frac{1}{r}$ and $\frac{1}{s}$. Because $x^2f(\frac{1}{x})=cx^2+bx+a$, by \href{https://artofproblemsolving.com/wiki/index.php?title=Vieta\%27s\_formulas}{Vieta's formulas}, $\frac{1}{r}+\frac{1}{s}=\boxed{\frac{-b}{c}}$, which is answer choice $\fbox{E}$.

 


	\item (\textit{AHSME 1965}) One side of a given triangle is 18 inches. Inside the triangle a line segment is drawn parallel to this side forming a trapezoid 
whose area is one-third of that of the triangle. The length of this segment, in inches, is: 


 $\textbf{(A)}\ 6\sqrt {6} \qquad  \textbf{(B) }\ 9\sqrt {2} \qquad  \textbf{(C) }\ 12 \qquad  \textbf{(D) }\ 6\sqrt{3}\qquad \textbf{(E) }\ 9$


 
 \textbf{Answer: }A



\textbf{Solution}


                    \begin{figure}[H]
                        centering
                        \includegraphics[width=0.5\textwidth]{Solution/AHSME_I_1965_8_Solution_Figure_1.png}
                    \end{figure}
                    

 Let the given triangle be $\triangle ABE$ with $AB=18$. Also, let $\overline{CD} \parallel \overline{AB}$ with $C$ on $\overline{BE}$ and $D$ on $\overline{AE}$. We know from the problem that $[ABCD]=\frac{1}{3}[\triangle ABE]$, so $[\triangle DCE]=\frac{2}{3}[\triangle ABE]$. Because $\overline{DC} \parallel \overline{AB}$, $\triangle DCE \sim \triangle ABE$ by \href{https://artofproblemsolving.com/wiki/index.php?title=AA\_similarity}{AA similarity}. Because the ratio of these two triangles' areas is $\frac{2}{3}$, the ratio of their sides must be $\sqrt{\frac{2}{3}}=\frac{\sqrt{6}}{3}$. Thus, $CD=\frac{\sqrt{6}}{3}AB=\frac{\sqrt{6}}{3}*18=\boxed{6\sqrt{6}}$, which is answer choice $\fbox{A}$.

 


	\item (\textit{AHSME 1965}) The vertex of the parabola $y = x^2 - 8x + c$ will be a point on the $x$-axis if the value of $c$ is: 


 $\textbf{(A)}\ - 16 \qquad  \textbf{(B) }\ - 4 \qquad  \textbf{(C) }\ 4 \qquad  \textbf{(D) }\ 8 \qquad  \textbf{(E) }\ 16$


 
 \textbf{Answer: }E



\textbf{Solution}

Notice that if the vertex of a parabola is on the x-axis, then the x-coordinate of the vertex must be a solution to the quadratic. Since the quadratic is strictly increasing on either side of the vertex, the solution must have double multiplicity, or the quadratic is a perfect square trinomial. This means that for the vertex of $y = x^2 - 8x + c$ to be on the x-axis, the trinomial must be a perfect square, and have discriminant of zero. So,

 
 \begin{align*}     0 &= b^2-4ac\\     0 &= (-8)^2-4c\\     c &=  64\\     c &= 16\\ \end{align*}


 Therefore $c=16$, and our answer is $\boxed{\textbf{(E)}}$.

 


	\item (\textit{AHSME 1965}) The statement $x^2 - x - 6 < 0$ is equivalent to the statement: 


 $\textbf{(A)}\ - 2 < x < 3 \qquad  \textbf{(B) }\ x > - 2 \qquad  \textbf{(C) }\ x < 3 \\ \textbf{(D) }\ x > 3 \text{ and }x < - 2 \qquad  \textbf{(E) }\ x > 3 \text{ or }x < - 2$


 
 \textbf{Answer: }A



\textbf{Solution}

To solve this problem, we may begin by factoring $x^2-x-6$ as $(x-3)(x+2)$. This is an upward opening parabola, therefore the solutions to $(x-3)(x+2) < 0$ are between the roots of the equation. That means our solutions are all $x$ such that $-2 < x < 3$, or simply $\boxed{\textbf{(A)}}$.

 


	\item (\textit{AHSME 1965}) Consider the statements: $\newline$
$I: (\sqrt{-4})(\sqrt {-16}) = \sqrt{(-4)(-16)}$, $\newline$
$II: \sqrt{(-4)(-16)} = \sqrt{64}$, $\newline$
$III: \sqrt{64} = 8$. $\newline$
Of these the following are incorrect. 


 $\textbf{(A)}\ \text{none} \qquad  \textbf{(B) }\ \text{I only} \qquad  \textbf{(C) }\ \text{II only} \qquad  \textbf{(D) }\ \text{III only}\qquad \textbf{(E) }\ \text{I and III only}$


 
 \textbf{Answer: }B



\textbf{Solution}

Statement I is incorrect, because $(\sqrt{-4})(\sqrt{-16})=(i\sqrt{4})(i\sqrt{16})=i^2(\sqrt{4})(\sqrt{16})=-\sqrt{64}$. $\newline$Statement II is correct, because $(-4)(-16)=64$, so $\sqrt{(-4)(-16)}=\sqrt{64}$. $\newline$Statement III is correct, because $8^2=64$ and $8 \geq 0$ (so it is the \href{https://artofproblemsolving.com/wiki/index.php?title=Square\_root}{principal square root} of 64). $\newline$Thus, only statement I is incorrect, so we choose answer $\fbox{B}$.

 


	\item (\textit{AHSME 1965}) A rhombus is inscribed in $\triangle ABC$ in such a way that one of its vertices is $A$ and two of its sides lie along $AB$ and $AC$. 
If $\overline{AC} = 6$ inches, $\overline{AB} = 12$ inches, and $\overline{BC} = 8$ inches, the side of the rhombus, in inches, is: 


 $\textbf{(A)}\ 2 \qquad  \textbf{(B) }\ 3 \qquad  \textbf{(C) }\ 3 \frac {1}{2} \qquad  \textbf{(D) }\ 4 \qquad  \textbf{(E) }\ 5$


 
 \textbf{Answer: }D



\textbf{Solution}


                    \begin{figure}[H]
                        centering
                        \includegraphics[width=0.5\textwidth]{Solution/AHSME_I_1965_12_Solution_Figure_1.png}
                    \end{figure}
                    

 As in the diagram, suppose the rhombus $ADEF$ is inscribed in $\triangle ABC$ with $D$ on $\overline{AC}$, $E$ on $\overline{BC}$, and $F$ on $\overline{AB}$. Let the side length of the rhombus be $x$. Because a rhombus is a parallelogram, its opposite sides are parallel. Thus, by \href{https://artofproblemsolving.com/wiki/index.php?title=AA\_similarity}{AA similarity}, $\triangle BFE \sim \triangle EDC$, and so $\frac{FE}{FB}=\frac{DC}{DE}$. Because $\overline{FE}$ and $\overline{DE}$ are sides of the rhombus, they both have length $x$. Furthermore, $FB=12-x$, and $DC=6-x$, because, combined with sides of the rhombus, they form sides of the triangle. Thus, by substituting into the proportion derived above, we see that:\begin{align*}\frac{FE}{FB}&=\frac{DC}{DE} \\\frac{x}{12-x}&=\frac{6-x}{x} \\x^2&=72-12x-6x+x^2 \\18x&=72 \\x&=4 \\\end{align*}


 Thus, the side of the rhombus has length $\boxed{\text{(D) }4}$.

 


	\item (\textit{AHSME 1965}) Let $n$ be the number of number-pairs $(x,y)$ which satisfy $5y - 3x = 15$ and $x^2 + y^2 \le 16$. Then $n$ is: 


 $\textbf{(A)}\ 0 \qquad  \textbf{(B) }\ 1 \qquad  \textbf{(C) }\ 2 \qquad  \textbf{(D) }\ \text{more than two, but finite}\qquad \textbf{(E) }\ \text{greater than any finite number}$


 
 \textbf{Answer: }E



\textbf{Solution}


                    \begin{figure}[H]
                        centering
                        \includegraphics[width=0.5\textwidth]{Solution/AHSME_I_1965_13_Solution_Figure_1.png}
                    \end{figure}
                    

 For an ordered pair $(x,y)$ to satisfy the restrictions, the point $(x,y)$ must lie on the graph of the line $5y-3x=15$ (shown in blue in the diagram) and on the closed \href{https://artofproblemsolving.com/wiki/index.php?title=Disk}{disk} given by $x^2+y^2 \leq 16$ (whose boundary is red in the diagram). Because the $y$-intercept of the line, $(0,3)$, is within the disk (and not on the boundary of the disk), the line cannot be tangent to the disk, and so the disk must contain infinitely many of the points on the line. Thus, our answer is $\fbox{\textbf{(E)} greater than any finite number}$

 


	\item (\textit{AHSME 1965}) The sum of the numerical coefficients in the complete expansion of $(x^2 - 2xy + y^2)^7$ is: 


 $\textbf{(A)}\ 0 \qquad  \textbf{(B) }\ 7 \qquad  \textbf{(C) }\ 14 \qquad  \textbf{(D) }\ 128 \qquad  \textbf{(E) }\ 128^2$


 
 \textbf{Answer: }A



\textbf{Solution}

Notice that the given equation, $(x^2 - 2xy + y^2)^7$ can be factored into $(x-y)^{14} = \sum_{k=0}^{14} \binom{14}{k}(-1)^k\cdot x^{14-k}\cdot y^k$. 

 Notice that if we plug in $x = y = 1$, then we are simply left with the sum of the coefficients from each term. Therefore, the sum of the coefficients is $(1-1)^7 = 0, \boxed{\textbf{(A)}}$.

 


	\item (\textit{AHSME 1965}) The symbol $25_b$ represents a two-digit number in the base $b$. If the number $52_b$ is double the number $25_b$, then $b$ is: 


 $\textbf{(A)}\ 7 \qquad  \textbf{(B) }\ 8 \qquad  \textbf{(C) }\ 9 \qquad  \textbf{(D) }\ 11 \qquad  \textbf{(E) }\ 12$


 
 \textbf{Answer: }B



\textbf{Solution}

We begin by noting that $25_b = 2b+5$ and $52_b = 5b+2$. The problem tells us that $2*25_b=52_b$, so $2(2b+5)=5b+2$. Solving for b yields the answer $\boxed{\textbf{(B) }8}$.

 


	\item (\textit{AHSME 1965}) Let line $AC$ be perpendicular to line $CE$. Connect $A$ to $D$, the midpoint of $CE$, and connect $E$ to $B$, 
the midpoint of $AC$. If $AD$ and $EB$ intersect in point $F$, and $\overline{BC} = \overline{CD} = 15$ inches, 
then the area of triangle $DFE$, in square inches, is: 


 $\textbf{(A)}\ 50 \qquad  \textbf{(B) }\ 50\sqrt {2} \qquad  \textbf{(C) }\ 75 \qquad  \textbf{(D) }\ \frac{15}{2}\sqrt{105}\qquad \textbf{(E) }\ 100$


 
 \textbf{Answer: }C



\textbf{Solution}


                    \begin{figure}[H]
                        centering
                        \includegraphics[width=0.5\textwidth]{Solution/AHSME_I_1965_16_Solution_Figure_1.png}
                    \end{figure}
                    

 Draw $\overline{AE}$, as seen in the diagram. From the problem, we know that $\overline{EB}$ and $\overline{AD}$ are \href{https://artofproblemsolving.com/wiki/index.php?title=Median\_of\_a\_triangle}{medians} of $\triangle ACE$. Let $G$ be the midpoint of $\overline{AE}$. Then, $\overline{CG}$ is also a median of $\triangle ACE$, and it goes through $\triangle ACE$'s \href{https://artofproblemsolving.com/wiki/index.php?title=Centroid}{centroid}, $F$. Because medians divide their triangle into $6$ smaller triangles of equal area, we know that $[\triangle DFE]=\frac{1}{6}[\triangle ACE]$. Because $[\triangle ACE]=\frac{1}{2}*(15+15)^2=\frac{900}{2}=450$, $[\triangle DFE]=\frac{450}{6}=75$. Thus, our answer is $\boxed{\textbf{(C) }75}$.

 


	\item (\textit{AHSME 1965}) Given the true statement: The picnic on Sunday will not be held only if the weather is not fair. We can then conclude that: 


 $\textbf{(A)}\ \text{If the picnic is held, Sunday's weather is undoubtedly fair.} \\ \textbf{(B) }\ \text{If the picnic is not held, Sunday's weather is possibly unfair.} \\ \textbf{(C) }\ \text{If it is not fair Sunday, the picnic will not be held.} \\ \textbf{(D) }\ \text{If it is fair Sunday, the picnic may be held.} \\ \textbf{(E) }\ \text{If it is fair Sunday, the picnic must be held.}$


 
 \textbf{Answer: }E



\textbf{Solution}

The given statement tells us that the weather being not fair is the only thing that can cancel the picnic on Sunday. It does not say that the weather being not fair necessarily cancels the picnic, so choice (A) is out. If the weather is fair, however, nothing can stop  the picnic from happening, so choice $\fbox{\textbf{(E)}}$ is correct. $\newline$Choice (B) is incorrect, because if the picnic is cancelled, then the weather \textit{must} be unfair, because the weather is the only thing that can stop the picnic. $\newline$Choice (C) is incorrect, because it is the \href{https://artofproblemsolving.com/wiki/index.php?title=Contrapositive}{contrapositive} of choice (A), so (A) and (C) are logically equivalent. $\newline$Choice (D) is incorrect, because the picnic \textit{must} be held if foul weather is not there to prevent it.

 


	\item (\textit{AHSME 1965}) If $1 - y$ is used as an approximation to the value of $\frac {1}{1 + y}, |y| < 1$, the ratio of the error made to the correct value is: 


 $\textbf{(A)}\ y \qquad  \textbf{(B) }\ y^2 \qquad  \textbf{(C) }\ \frac {1}{1 + y} \qquad  \textbf{(D) }\ \frac{y}{1+y}\qquad \textbf{(E) }\ \frac{y^2}{1+y}\qquad$


 
 \textbf{Answer: }B



\textbf{Solution}

The error made in this approximation is $\frac{1}{1+y}-(1-y)$, and the correct value is $\frac{1}{1+y}$. Taking the ratio of these two values, we have:\begin{align*} \\\frac{\frac{1}{1+y}-(1-y)}{\frac{1}{1+y}}&=\frac{[\frac{1}{1+y}-(1-y)][1+y]}{[\frac{1}{1+y}][1+y]} \\&=\frac{1-(1-y^2)}{1} \\&=y^2 \\\end{align*}


 Thus, our answer is $\boxed{\textbf{(B) }y^2}$

 


	\item (\textit{AHSME 1965}) If $x^4 + 4x^3 + 6px^2 + 4qx + r$ is exactly divisible by $x^3 + 3x^2 + 9x + 3$, the value of $(p + q)r$ is: 


 $\textbf{(A)}\ - 18 \qquad  \textbf{(B) }\ 12 \qquad  \textbf{(C) }\ 15 \qquad  \textbf{(D) }\ 27 \qquad  \textbf{(E) }\ 45 \qquad$


 
 \textbf{Answer: }C



\textbf{Solution 1}

Let $f(x)=x^3+3x^2+9x+3$ and $g(x)=x^4+4x^3+6px^2+4qx+r$. 

 Let 3 roots of $f(x)$ be $r_1, r_2$ and $r_3$. As  $f(x)|g(x)$ , 3 roots of 4 roots of $g(x)$ will be same as roots of $f(x)$. Let  the 4th root of $g(x)$ be $r_4$.  By \href{https://artofproblemsolving.com/wiki/index.php?title=Vieta\%27s\_Formulas}{Vieta's Formulas}

 In $f(x)$

 $r_1+r_2+r_3=-3$

 $r_1r_2+r_2r_3+r_1r_3=9$

 $r_1r_2r_3=-3$

 In $g(x)$

 $r_1+r_2+r_3+r_4=-4$

 $=>r_4=-1$

 $r_1r_2+r_1r_3+r_1r_4+r_2r_3+r_2r_4+r_3r_4=6p$

 $=>r_1r_2+r_1r_3+r_2r_3+r_4(r_1+r_2+r_3)=6p$

 $=>9+(-1)(-3)=6p$

 $=>p=2$

 $r_1r_2r_3+r_1r_3r_4+r_1r_2r_4+r_2r_3r_4=-4q$

 $=>-3+r_4(r_1r_3+r_1r_2+r_2r_3)=-4q$

 $=>-3+(-1)(9)=-4q$

 $=>q=3$

 $r_1r_2r_3r_4=r$

 $=>(-3)(-1)=r$

 $=>r=3$

 so $(p+q)r=\boxed{\textbf{(C) }15}$

 By ~Ahmed\_Ashhab

 



\textbf{Solution 2}

Notice that to obtain the $x^4$ term one must multiply $x^4+4x^3+6px^2+4qx+r$ by some linear function of the form $x-a$. Looking at the $x^3$ term, it is clear that $a$ must equal $1$. Therefore by multiplying $x^4+4x^3+6px^2+4qx+r$ by $x+1$, the product will be $x^4+4x^3+12x^2+12x+3$. Therefore  $p=2$, $q=3$, $r=3$. Thus $(2+3)3=\boxed{\textbf{(C) }15}$

 


	\item (\textit{AHSME 1965}) For every $n$ the sum of n terms of an arithmetic progression is $2n + 3n^2$. The $r$th term is: 


 $\textbf{(A)}\ 3r^2 \qquad  \textbf{(B) }\ 3r^2 + 2r \qquad  \textbf{(C) }\ 6r - 1 \qquad  \textbf{(D) }\ 5r + 5 \qquad  \textbf{(E) }\ 6r+2\qquad$


 
 \textbf{Answer: }C



\textbf{Solution}

In order to calculate the $r$th term of this arithmetic sequence, we can subtract the sum of the first $r-1$ terms from the sum of the first $r$ terms of the sequence. Plugging in $r$ and $r-1$ as values of $n$ in the given expression and subtracting yields 
 \[(3r^2+2r)-(3r^2-4r+1).\] Simplifying gives us the final answer of $\boxed{\textbf{(C) }6r-1}$.

 


	\item (\textit{AHSME 1965}) It is possible to choose $x > \frac {2}{3}$ in such a way that the value of $\log_{10}(x^2 + 3) - 2 \log_{10}x$ is 


 $\textbf{(A)}\ \text{negative} \qquad  \textbf{(B) }\ \text{zero} \qquad  \textbf{(C) }\ \text{one} \\ \textbf{(D) }\ \text{smaller than any positive number that might be specified} \\ \textbf{(E) }\ \text{greater than any positive number that might be specified}$


 
 \textbf{Answer: }D



\textbf{Solution}

By the rules of \href{https://artofproblemsolving.com/wiki/index.php?title=Logarithms}{logarithms}, $\log_{10}(x^2+3)-2\log_{10} x=\log_{10}(\frac{x^2+3}{x^2})=\log_{10}(1+\frac{3}{x^2})$. As $x$ goes to infinity, $1+\frac{3}{x^2}$ gets arbitrarily close to $1$ (without ever reaching it), so $\log_{10}(1+\frac{3}{x^2})$ gets arbitrarily close to $\log_{10}(1)=0$ (without ever reaching it). Furthermore, because $1+\frac3{x^2} > 1$, $\log(1+\frac3{x^2})$ is never negative. Thus, we can choose a real $x>\frac{2}{3}$ such that the given expression is $\fbox{\textbf{(D) }smaller than any positive number that might be specified}$.

 


	\item (\textit{AHSME 1965}) If $a_2 \neq 0$ and $r$ and $s$ are the roots of $a_0 + a_1x + a_2x^2 = 0$, then the equality 
$a_0 + a_1x + a_2x^2 = a_0\left (1 - \frac {x}{r} \right ) \left (1 - \frac {x}{s} \right )$ holds: 


 $\textbf{(A)}\ \text{for all values of }x, a_0\neq 0 \qquad \textbf{(B) }\ \text{for all values of }x \\ \textbf{(C) }\ \text{only when }x = 0 \qquad \textbf{(D) }\ \text{only when }x = r \text{ or }x = s \\ \textbf{(E) }\ \text{only when }x = r \text{ or }x = s, a_0 \neq 0$


 
 \textbf{Answer: }A



\textbf{Solution}

$a_0$ cannot be $0$, because $a_2 \neq 0$, so the polynomial can take non-zero values when $a_0=0$ and thereby not satisfy the equation. Expanding $a_0(1-\frac{x}{r})(1-\frac{x}{s})$ gives us $a_0(1-\frac{x}{r}-\frac{x}{s}+\frac{x^2}{rs})$. Using \href{https://artofproblemsolving.com/wiki/index.php?title=Vieta\%27s\_formulas}{Vieta's formulas}, we see that $rs=\frac{a_0}{a_2}$, so $a_0=a_2rs$. Plugging this into the expanded right hand side of the given equation, we see that that side equals $a_0-a_2(r+s)x+a_2x^2$. Using Vieta's formulas again, we equate this expression to $a_0+a_1x+a_2x^2$, which is the left hand side of the given equation. Thus, as long as $a_0 \neq 0$, the equation holds for all values of $x$. This fact corresponds with answer choice $\fbox{\textbf{(A)}}$.

 


	\item (\textit{AHSME 1965}) If we write $|x^2 - 4| < N$ for all $x$ such that $|x - 2| < 0.01$, the smallest value we can use for $N$ is: 


 $\textbf{(A)}\ .0301 \qquad  \textbf{(B) }\ .0349 \qquad  \textbf{(C) }\ .0399 \qquad  \textbf{(D) }\ .0401 \qquad  \textbf{(E) }\ .0499\qquad$


 
 \textbf{Answer: }D



\textbf{Solution}

There are two cases: (i) $x-2>0$ and (ii) $x-2<0$. In case (i), $0<x-2<0.01$, so $4<x+2<4.01$. In case (ii), $-0.01<x-2<0$, so $3.99<x+2<4$. Because $(x+2)$ in case (i) can take a larger value, we use it to determine the upper bound for $|x^2-4|$ under the given restrictions. Now, we can see that $|x^2-4|=|x-2||x+2|<0.01*4.01=\boxed{0.0401}$, which is answer choice $\fbox{\textbf{(D)}}$.

 


	\item (\textit{AHSME 1965}) Given the sequence $10^{\frac {1}{11}},10^{\frac {2}{11}},10^{\frac {3}{11}},\ldots,10^{\frac {n}{11}}$, 
the smallest value of n such that the product of the first $n$ members of this sequence exceeds $100000$ is: 


 $\textbf{(A)}\ 7 \qquad  \textbf{(B) }\ 8 \qquad  \textbf{(C) }\ 9 \qquad  \textbf{(D) }\ 10 \qquad  \textbf{(E) }\ 11$


 
 \textbf{Answer: }E



\textbf{Solution}

Note that the given sequence is a \href{https://artofproblemsolving.com/wiki/index.php?title=Geometric\_sequence}{geometric sequence} with a common ratio $10^{\frac{1}{11}}$. Let the product of the first $n$ terms of the sequence be denoted $P_n$. It is a consequence of the laws of exponents that $P_1=(10^{\frac{1}{11}})^1$, $P_2=(10^{\frac{1}{11}})^{1+2}$, and, in general, $P_n=(10^{\frac{1}{11}})^{T_n}$, where $T_n$ denotes the $n$th \href{https://artofproblemsolving.com/wiki/index.php?title=Triangular\_number}{triangular number}. Setting $P_n$ equal to $100,000$, we see that:\begin{align*} \\(10^{\frac{1}{11}})^{\frac{n(n+1)}{2}}&=10^5 \\\frac{1}{11}*\frac{n(n+1)}{2}&=5 \\n(n+1)&=110 \\n^2+n-110&=0 \\(n-10)(n+11)&=0 \\\end{align*}
Because $n$ must be positive, we are left with $n=10$. Given this information, choice (D) may seem appealing. Do not be fooled. The product of the first $n$ terms is \textit{exactly} $100,000$, but the problem asks for the smallest $n$ such that $P_n$ \textit{exceeds} 100,000. Thus, the minimum value of $n$ which satsifies the problem is $\boxed{\textbf{(E) }11}$.

 


	\item (\textit{AHSME 1965}) Let $ABCD$ be a quadrilateral with $AB$ extended to $E$ so that $\overline{AB} = \overline{BE}$. 
Lines $AC$ and $CE$ are drawn to form $\angle{ACE}$. For this angle to be a right angle it is necessary that quadrilateral $ABCD$ have: 


 $\textbf{(A)}\ \text{all angles equal} \qquad \textbf{(B) }\ \text{all sides equal} \\ \textbf{(C) }\ \text{two pairs of equal sides} \qquad \textbf{(D) }\ \text{one pair of equal sides} \\ \textbf{(E) }\ \text{one pair of equal angles}$


 
 \textbf{Answer: }D



\textbf{Solution}


                    \begin{figure}[H]
                        centering
                        \includegraphics[width=0.5\textwidth]{Solution/AHSME_I_1965_25_Solution_Figure_1.png}
                    \end{figure}
                    

 Because $\triangle ACE$ is right, the midpoint of its hypoteneuse (namely, $B$) is its \href{https://artofproblemsolving.com/wiki/index.php?title=Orthocenter}{orthocenter}. Thus, $AB=BC$, and so two side lengths of quadrilateral $ABCD$ are equal. The placement of $D$ is irrelevant. Thus, our answer is $\fbox{\textbf{(D) }one pair of equal sides}$.

 


	\item (\textit{AHSME 1965}) For the numbers $a, b, c, d, e$ define $m$ to be the arithmetic mean of all five numbers; 
$k$ to be the arithmetic mean of $a$ and $b$; $l$ to be the arithmetic mean of $c, d$, and $e$; 
and $p$ to be the arithmetic mean of $k$ and $l$. Then, no matter how $a, b, c, d$, and $e$ are chosen, we shall always have: 


 $\textbf{(A)}\ m = p \qquad  \textbf{(B) }\ m \ge p \qquad  \textbf{(C) }\ m > p \qquad \\ \textbf{(D) }\ m < p\qquad \textbf{(E) }\ \text{none of these}$


 
 \textbf{Answer: }E



\textbf{Solution}

We shall begin by eliminating some options through counterexamples. If $a=b=c=d=e$, then $m=k=l=a$, so $p=a$, and $m=p$. Answers (C) and (D) do not allow for $m=p$, so they can be eliminated. If we set $a=60$ and $b=c=d=e=0$, then $m=12$, $k=30$, $l=0$, and $p=15$. Here, $m<p$, so we can throw out options (A) and (B) as well. Now, we are left with only option $\fbox{\textbf{(E) }none of these}$.

 


	\item (\textit{AHSME 1965}) When $y^2 + my + 2$ is divided by $y - 1$ the quotient is $f(y)$ and the remainder is $R_1$. 
When $y^2 + my + 2$ is divided by $y + 1$ the quotient is $g(y)$ and the remainder is $R_2$. If $R_1 = R_2$ then $m$ is: 


 $\textbf{(A)}\ 0 \qquad  \textbf{(B) }\ 1 \qquad  \textbf{(C) }\ 2 \qquad  \textbf{(D) }\ - 1 \qquad  \textbf{(E) }\ \text{an undetermined constant}$


 
 \textbf{Answer: }A



\textbf{Solution}

Let $h(y)=y^2+my+2$

 $h(y)=y^2+my+2=f(y)(y-1)+R_1$

 $h(y)=y^2+my+2=g(y)(y+1)+R_2$

 $h(1)=3+m=R_1$

 $h(-1)=3-m=R_2$

 $m=0$

 Thus, the answer is $\boxed{\textbf{(A) }0}$.

 


 



\textbf{Solution 2}

As in the other solution, let $h(y)=y^2+my+2$. By dividing $h(y)$ by $(y-1)$ (through either \href{https://artofproblemsolving.com/wiki/index.php?title=Polynomial\_long\_division&action=edit&redlink=1}{long} or \href{https://artofproblemsolving.com/wiki/index.php?title=Synthetic\_division}{synthetic} division), we get a remainder $R_1=m+3$. Similarly, dividing $h(y)$ by $(y+1)$ yields the remainder $R_2=3-m$. Setting $R_1=R_2$, we see that $m+3=3-m$, which only holds when $m=0$, which corresponds to answer choice $\fbox{\textbf{(A)}}$.

 


	\item (\textit{AHSME 1965}) An escalator (moving staircase) of $n$ uniform steps visible at all times descends at constant speed. 
Two boys, $A$ and $Z$, walk down the escalator steadily as it moves, A negotiating twice as many escalator 
steps per minute as $Z$. $A$ reaches the bottom after taking $27$ steps while $Z$ reaches the bottom after taking $18$ steps. Then $n$ is: 


 $\textbf{(A)}\ 63 \qquad  \textbf{(B) }\ 54 \qquad  \textbf{(C) }\ 45 \qquad  \textbf{(D) }\ 36 \qquad  \textbf{(E) }\ 30$


 
 \textbf{Answer: }B



\textbf{Solution}

If we let $Z$'s speed be $z$ steps/minute, then $A$'s speed is $2z$ steps/minute. Let $t_a$ be the time $A$ spent on the escalator, and let $t_z$ be the time $Z$ spent on the escalator. Then, we know that $A$ walked down $2zt_a=27$ steps, and $Z$ walked down $zt_z=18$ steps. Dividing the first equation by the second, we see that:\begin{align*} \\\frac{2zt_a}{zt_z}&=\frac{27}{18} \\\frac{2t_a}{t_z}&=\frac{3}{2} \\t_a&=\frac{3}{4}t_z \\\end{align*}


 Thus, because $A$ was on the escalator for $\frac{3}{4}$ as long as $Z$ was, $A$ only gained $\frac{3}{4}$ as many "free" steps (i.e. steps that do not have to be taken because the escalator is moving down). We know that $A$ gained $(n-27)$ free steps, and $Z$ gained $(n-18)$ free steps. Thus we have the following equation: $n-27=\frac{3}{4}(n-18)$. Solving for $n$ gives us $\boxed{\textbf{(B) }54}$.

 


	\item (\textit{AHSME 1965}) Of $28$ students taking at least one subject the number taking Mathematics and English only equals the number 
taking Mathematics only. No student takes English only or History only, and six students take Mathematics and 
History, but not English. The number taking English and History only is five times the number taking all three subjects. 
If the number taking all three subjects is even and non-zero, the number taking English and Mathematics only is: 


 $\textbf{(A)}\ 5 \qquad  \textbf{(B) }\ 6 \qquad  \textbf{(C) }\ 7 \qquad  \textbf{(D) }\ 8 \qquad  \textbf{(E) }\ 9$


 
 \textbf{Answer: }A



\textbf{Solution}

The number of students taking all three subjects cannot be greater than or equal to $4$, because then there would be $20$ or more students taking English and History only, and we know that $6$ students are taking Mathematics and History only. These three values sum to a number greater than $28$, the total number of students. Thus, because we know that the number of students taking all three subjects is a non-zero even number, there must be $2$ of them. Thus, $10$ people must be taking only English and History. With the knowledge of $6$ students taking only Mathematics and History, we know the classes of $18$ students so far, so we have $10$ students left who are either taking Mathematics and English only or Mathematics only. Because these two categories have an equal number of students, the number of students taking English and Mathematics only is $\boxed{\textbf{(A) }5}$.

 


	\item (\textit{AHSME 1965}) Let $BC$ of right triangle $ABC$ be the diameter of a circle intersecting hypotenuse $AB$ in $D$. 
At $D$ a tangent is drawn cutting leg $CA$ in $F$. This information is not sufficient to prove that 


 $\textbf{(A)}\ DF \text{ bisects }CA \qquad  \textbf{(B) }\ DF \text{ bisects }\angle CDA \\ \textbf{(C) }\ DF = FA \qquad  \textbf{(D) }\ \angle A = \angle BCD \qquad  \textbf{(E) }\ \angle CFD = 2\angle A$


 
 \textbf{Answer: }B



\textbf{Solution 1}


                    \begin{figure}[H]
                        centering
                        \includegraphics[width=0.5\textwidth]{Solution/AHSME_I_1965_30_Solution_Figure_1.png}
                    \end{figure}
                    

 We will prove every result except for $\fbox{B}$. 

 By \href{https://artofproblemsolving.com/wiki/index.php?title=Thales\%27\_Theorem}{Thales' Theorem}, $\angle CDB=90^\circ$ and so $\angle CDA= 90^\circ$. $FC$ and $FD$ are both tangents to the same circle, and hence equal. Let $\angle CFD=\alpha$. Then $\angle FDC = \frac{180^\circ - \alpha}{2}$, and so $\angle FDA = \frac{\alpha}{2}$. We also have $\angle AFD = 180^\circ - \alpha$, which implies $\angle FAD=\frac{\alpha}{2}$. This means that $CF=DF=FA$, so $DF$ indeed bisects $CA$. We also know that $\angle BCD=90-\frac{180^\circ - \alpha}{2}=\frac{\alpha}{2}$, hence $\angle A = \angle BCD$. And $\angle CFD=2\angle A$ as $\alpha = \frac{\alpha}{2}\times 2$. 

 Since all of the results except for $B$ are true, our answer is $\fbox{B}$.

 



\textbf{Solution 2}

It's easy to verify that $\angle CDA$ always equals $90^\circ$. Since $\angle CDF$ changes depending on the sidelengths of the triangle, we cannot be certain that $\angle CDF=45^\circ$. Hence our answer is $\fbox{B}$.

 


	\item (\textit{AHSME 1965}) The number of real values of $x$ satisfying the equality $(\log_ax)(\log_bx) = \log_ab$, where $a > 0, b > 0, a \neq 1, b \neq 1$, is: 


 $\textbf{(A)}\ 0 \qquad  \textbf{(B) }\ 1 \qquad  \textbf{(C) }\ 2 \qquad  \textbf{(D) }\ \text{a finite integer greater than 2}\qquad \textbf{(E) }\ \text{not finite}$


 
 \textbf{Answer: }C



\textbf{Solution}

From the properties of \href{https://artofproblemsolving.com/wiki/index.php?title=Logarithms}{logarithms}, we know that $(\log_ax)(\log_xb)=\log_ab$, and $\log_xb=\frac{1}{\log_bx}$. Thus, to satisfy the given equation, we are looking for values of $x$ when $\log_bx=\frac{1}{\log_bx}$, or when $\log_bx=\pm 1$. This happens either when $x=b$ or $x=\frac{1}{b}$. Thus we have $2$ solutions, consistent with answer choice $\fbox{\textbf{(C)}}$.

 


	\item (\textit{AHSME 1965}) An article costing $C$ dollars is sold for $100 at a loss of $x$ percent of the selling price. 
It is then resold at a profit of $x$ percent of the new selling price $S'$. 
If the difference between $S'$ and $C$ is $1\frac {1}{9}$ dollars, then $x$ is: 


 $\textbf{(A)}\ \text{undetermined} \qquad  \textbf{(B) }\ \frac {80}{9} \qquad  \textbf{(C) }\ 10 \qquad  \textbf{(D) }\ \frac{95}{9}\qquad \textbf{(E) }\ \frac{100}{9}$


 
 \textbf{Answer: }C



\textbf{Solution}

The magnitude of the loss after the first sale is $C-100$, which equals $x$\% of the selling price, $\$100$. Thus, $C-100=100*\frac{x}{100}$, and so $C=x+100$. The profit made after the second sale, $S'-100$, is $x$% of the new selling price, and this quantity is represented by $S'*\frac{x}{100}$. Equating these two expressions, we see that $S'-100=S'*\frac{x}{100}$, and so $S'=\frac{10,000}{100-x}$. Because we know that the difference between $S'$ and $C$ is $\frac{10}{9}$, we can solve the following equation:\begin{align*} \\\frac{10,000}{100-x}-(x+100)&=\frac{10}{9} \\10,000-(100+x)(100-x)&=\frac{10}{9}(100-x) \\10,000-(10,000-x^2)&=\frac{1000}{9}-\frac{10}{9}x \\x^2+\frac{10}{9}x-\frac{1000}{9}&=0 \\9x^2+10x-1000&=0 \\(9x+100)(x-10)&=0 \\\end{align*}
Because $x \geq 0$, $x=10$, and so we choose answer $\boxed{\textbf{(C) }10}$.

 


	\item (\textit{AHSME 1965}) If the number $15!$, that is, $15 \cdot 14 \cdot 13 \dots 1$, ends with $k$ zeros when given to the base $12$ and ends with $h$ zeros 
when given to the base $10$, then $k + h$ equals: 


 $\textbf{(A)}\ 5 \qquad  \textbf{(B) }\ 6 \qquad  \textbf{(C) }\ 7 \qquad  \textbf{(D) }\ 8 \qquad  \textbf{(E) }\ 9$


 
 \textbf{Answer: }D



\textbf{Solution}

We can use \href{https://artofproblemsolving.com/wiki/index.php?title=Legendre\%27s\_Formula}{Legendre's Formula} to find the number of $0$s in base $10$
 \[\lfloor \frac{15}{5} \rfloor + \lfloor \frac{15}{25} \rfloor = 3\]So $h = 3$.Likewise, we are looking for the number of $2^2$s and $3$s that divide $15!$, so we use Legendre's again.
 \[\lfloor \frac{15}{2} \rfloor + \lfloor \frac{15}{4} \rfloor + \lfloor \frac{15}{8} \rfloor = 7 + 3 + 1 = 11\]
 \[\lfloor \frac{15}{3} \rfloor + \lfloor \frac{15}{9} \rfloor = 5 + 1 = 6\]Thus, $3^6 \vert 15!$ and $2^{11} \vert 15! \Rrightarrow (2^2)^5 \vert 15!$So $k = 5$, and $5+3 = 8$, which corresponds to answer $\fbox{\textbf{(D)}}$.

 


	\item (\textit{AHSME 1965}) For $x \ge 0$ the smallest value of $\frac {4x^2 + 8x + 13}{6(1 + x)}$ is: 


 $\textbf{(A)}\ 1 \qquad  \textbf{(B) }\ 2 \qquad  \textbf{(C) }\ \frac {25}{12} \qquad  \textbf{(D) }\ \frac{13}{6}\qquad \textbf{(E) }\ \frac{34}{5}$


 
 \textbf{Answer: }B



\textbf{Solution 1}

To begin, lets denote the equation, $\frac {4x^2 + 8x + 13}{6(1 + x)}$ as $f(x)$. Let's notice that:

 
 \begin{align*} f(x) & = \frac {4x^2 + 8x + 13}{6(1 + x)}\\\\ 	 & = \frac{4(x^2+2x) + 13}{6(x+1)}\\\\      & = \frac{4(x^2+2x+1-1)+13}{6(x+1)}\\\\      & = \frac{4(x+1)^2+9}{6(x+1)}\\\\      & = \frac{4(x+1)^2}{6(x+1)} + \frac{9}{6(1+x)}\\\\      & = \frac{2(x+1)}{3} + \frac{3}{2(x+1)}\\\\ \end{align*}


 After this simplification, we may notice that we may use calculus, or the \href{https://artofproblemsolving.com/wiki/index.php?title=AM-GM\_inequality}{AM-GM inequality} to finish this problem because $x\ge 0$, which implies that both $\frac{2(x+1)}{3} \text{ and } \frac{3}{2(x+1)}$ are greater than zero. Continuing with AM-GM:

 
 \begin{align*} \frac{\frac{2(x+1)}{3} + \frac{3}{2(x+1)}}{2} &\ge {\small \sqrt{\frac{2(x+1)}{3}\cdot \frac{3}{2(x+1)}}}\\\\ f(x) &\ge 2 \end{align*}


 Therefore, $f(x) = \frac {4x^2 + 8x + 13}{6(1 + x)} \ge 2$, $\boxed{\textbf{(B)}}$ 

 
$(\text{With equality when } \frac{2(x+1)}{3} = \frac{3}{2(x+1)}, \text{ or } x=\frac{1}{2})$

 



\textbf{Solution 2 (Calculus)}

Let $f(x)=\frac{4x^2 + 8x + 13}{6(1 + x)}$.Take the \href{https://artofproblemsolving.com/wiki/index.php?title=Derivative}{derivative} of $f(x)$ using the \href{https://artofproblemsolving.com/wiki/index.php?title=Quotient\_rule}{quotient rule}.\begin{align*} f(x) &= \frac{4x^2 + 8x + 13}{6(1 + x)} \\f'(x) &= \frac{1}{6}*\frac{(4x^2 + 8x + 13)'(1 + x) - (4x^2 + 8x + 13)(1 + x)'}{(1 + x)^2} \\&= \frac{(8x + 8)(1 + x) - (4x^2 + 8x + 13)(1)}{6(1 + x)^2} \\&= \frac{4x^2 + 8x - 5}{6(1 + x)^2} \\\end{align*}
Next, set the numerator equal to zero to find the $x$-value of the minimum:\begin{align*}4x^2+8x-5 &= 0 \\ (2x+5)(2x-1) &= 0 \\\end{align*}
From the problem, we know that $x \geq 0$, so we are left with $x=\frac{1}{2}$. Plugging $x=\frac{1}{2}$ into $f(x)$, we get:\begin{align*}f(\frac{1}{2})&=\frac{4(\frac{1}{2})^2+8(\frac{1}{2})+13}{6(1+(\frac{1}{2}))} \\&=\frac{1+4+13}{6(\frac{3}{2})} \\&=\frac{18}{9} \\&=2 \\\end{align*}


 Thus, our answer is $\boxed{\textbf{(B) }2}$.

 



\textbf{Solution 3 (answer choices, no AM-GM or calculus)}

We go from A through E and we look to find the smallest value so that $x \ge 0$, so we start from A:

 
 \[\frac{4x^2 + 8x + 13}{6x+6} = 1\]

 
 \[4x^2 + 8x + 13 = 6x + 6\]

 
 \[4x^2 + 2x + 7 = 0\]

 However by the quadratic formula there are no real solutions of $x$, so $x$ cannot be greater than 0. We move on to B:

 
 \[\frac{4x^2 + 8x + 13}{6x + 6} = 2\]

 
 \[4x^2 + 8x + 13 = 12x + 12\]

 
 \[4x^2 - 4x + 1 = 0\]

 
 \[(2x-1)^2 = 0\]

 There is one solution: $x = \frac{1}{2}$, which is greater than 0, so 2 works as a value. Since all the other options are bigger than 2 or invalid, the answer must be $\boxed{\textbf{B}}$

 


	\item (\textit{AHSME 1965}) The length of a rectangle is $5$ inches and its width is less than $4$ inches. The rectangle is folded so that two 
diagonally opposite vertices coincide. If the length of the crease is $\sqrt {6}$, then the width is: 


 $\textbf{(A)}\ \sqrt {2} \qquad  \textbf{(B) }\ \sqrt {3} \qquad  \textbf{(C) }\ 2 \qquad  \textbf{(D) }\ \sqrt{5}\qquad \textbf{(E) }\ \sqrt{\frac{11}{2}}$


 
 \textbf{Answer: }D



\textbf{Solution}


                    \begin{figure}[H]
                        centering
                        \includegraphics[width=0.5\textwidth]{Solution/AHSME_I_1965_35_Solution_Figure_1.png}
                    \end{figure}
                    

 Let the rectangle be $ABCD$ with $AB=CD=w$ and $AD=BC=5$, as in the diagram. We desire a line such that reflecting point $C$ across that line yields point $A$. For this to happen, the line must be perpendicular to the diagonal $\overline{AC}$, and it must go through the midpoint of $\overline{AC}$ (let it be point $M$). Let the intersection of this line with $\overline{BC}$ be point $X$ and with $\overline{AD}$ be point $Y$. From the problem, we know that $XY=\sqrt{6}$. By \href{https://artofproblemsolving.com/wiki/index.php?title=Congruent\_(geometry)\#HL\_Congruence}{HL congruence}, $\triangle AMY \cong \triangle CMX$, so $AM=CM=x$, where $x$ is some number. Furthermore, $XM=MY=\frac{\sqrt{6}}{2}$. By \href{https://artofproblemsolving.com/wiki/index.php?title=AA\_similarity}{AA similarity}, $\triangle MXC \sim \triangle BAC$, so $\frac{MX}{MC}=\frac{BA}{BC}$. $MX=\frac{\sqrt{6}}{2}$,  $MC=x$, $BA=w$, and $BC=5$, so we can rewrite this proportion to solve for $x$ in terms of $w$:\begin{align*}\frac{\sqrt{6}/2}{x}&=\frac{w}{5} \\\frac{5\sqrt{6}}{2}&=xw \\x&=\frac{5\sqrt{6}}{2w}\end{align*}
By the \href{https://artofproblemsolving.com/wiki/index.php?title=Pythagorean\_Theorem}{Pythagorean Theorem} on $\triangle ABC$, we know that $w^2+25=4x^2$, and we can plug in our new expression for $x$ into this equation to solve for $w$:\begin{align*}w^2+25&=4(\frac{5\sqrt{6}}{2w})^2 \\w^2+25&=\frac{25*6}{w^2} \\w^4+25w^2-150&=0 \\(w^2-5)(w^2+30)&=0\end{align*}
Because $w>0$, $w^2=5$, and so $w=\boxed{\sqrt{5}}$, which corresponds to answer choice $\fbox{\textbf{(D)}}$.

 


	\item (\textit{AHSME 1965}) Given distinct straight lines $OA$ and $OB$. From a point in $OA$ a perpendicular is drawn to $OB$; 
from the foot of this perpendicular a line is drawn perpendicular to $OA$. 
From the foot of this second perpendicular a line is drawn perpendicular to $OB$; 
and so on indefinitely. The lengths of the first and second perpendiculars are $a$ and $b$, respectively. 
Then the sum of the lengths of the perpendiculars approaches a limit as the number of perpendiculars grows beyond all bounds. This limit is: 


 $\textbf{(A)}\ \frac {b}{a - b} \qquad  \textbf{(B) }\ \frac {a}{a - b} \qquad  \textbf{(C) }\ \frac {ab}{a - b} \qquad  \textbf{(D) }\ \frac{b^2}{a-b}\qquad \textbf{(E) }\ \frac{a^2}{a-b}$


 
 \textbf{Answer: }E



\textbf{Solution 1}


                    \begin{figure}[H]
                        centering
                        \includegraphics[width=0.5\textwidth]{Solution/AHSME_I_1965_36_Solution_Figure_1.png}
                    \end{figure}
                    

 For simplicity, let the first perpendicular from $\overleftrightarrow{OA}$ to $\overleftrightarrow{OB}$ be $\overline{AB}$, and let the second perpendicular have foot $C$ on $\overleftrightarrow{OA}$. Further, let the perpendicular from $C$ to $\overleftrightarrow{OB}$ have foot $D$ and length $c$, as in the diagram. Also, let $\measuredangle OAB=\alpha$. From the problem, we have $AB=a$ and $BC=b$. By \href{https://artofproblemsolving.com/wiki/index.php?title=AA\_similarity}{AA similarity}, we have $\triangle OCB \sim \triangle OBA$, so $\measuredangle CBO=\alpha$ as well. In $\triangle ABC$, we see that $\sin\alpha=\frac{b}{a}$, and, in $\triangle CDB$, $\sin\alpha=\frac{c}{b}$. Equating these two expressions for $\sin\alpha$, we get that $\frac{b}{a}=\frac{c}{b}$, or, because $a,b,c>0$, $b=\sqrt{ac}$. Thus, $b$ is the \href{https://artofproblemsolving.com/wiki/index.php?title=Geometric\_mean}{geometric mean} of $a$ and $c$. Note that if we remove the first perpendicular (i.e. the one with length $a$), we are left with a smaller version of the original problem, which will have the same equation for the limit (but this time expressed in terms of $b$ and $c$ rather than $a$ and $b$). Thus, if we let the length of the fourth perpendicular be $d$, then $c$ will equal the geometric mean of $b$ and $d$, and so on for the infinitude of perpendiculars. Thus, because the length of a given perpendicular (except the first one) is the geometric mean of the two adjacent perpendiculars, the lengths of the perpendiculars form a \href{https://artofproblemsolving.com/wiki/index.php?title=Geometric\_sequence}{geometric sequence}. Because the sequence's first two terms are $a$ and $b$, it has common ratio $\frac{b}{a}$. Because $b<a$, the common ratio is positive and less than $1$, so the sequence's infinite \href{https://artofproblemsolving.com/wiki/index.php?title=Geometric\_sequence\#Sum}{geometric series} converges. This infinite sum is given by $\frac{a}{1-\frac{b}{a}}=\boxed{\frac{a^2}{a-b}}$, which is answer choice $\fbox{\textbf{(E)}}$.

 



\textbf{Solution 2 (Answer choices, intuition)}

Let $a$ and $b$ be measured with some units of length (say, meters). The limit of the sum of the lengths of the perpendiculars would, then, be measured in meters as well. Performing \href{https://artofproblemsolving.com/wiki/index.php?title=Dimensional\_analysis&action=edit&redlink=1}{dimensional analysis} on each of the answer choices, we can eliminate options (A) and (B), because they have units $\frac{\text{meters}}{\text{meters}}$, so they are dimensionless. Unfortunately, our other three options all have units of meters. Now, we have a $\frac{1}{3}$ chance of guessing the answer correctly, but we can go further. Think of what happens as $\measuredangle BOA$ approaches $0^{\circ}$. Then, $a$ approaches $0$, and the perpendiculars have nearly no space to "bounce between" the two lines, and so they \textit{likely} have zero total length. Likewise, if we think of what happens as $\measuredangle BOA$ approaches $90^{\circ}$, all of the perpendiculars except the first one (with length $a$) go to zero. From this intuition, one would think that $a$ would have a larger impact upon the total length of the perpendiculars. Out of the three answer choices left, this conjecture is only consistent with choice $\boxed{\textbf{(E) }\frac{a^2}{a-b}}$. To reinforce this decision, think about what happens as $\measuredangle BOA$ approaches $90^{\circ}$ again. We would \textit{expect} the sum to approach $a$ and $b$ to approach $0$. Plugging in $b=0$ into the expression in choice (E), we get $\frac{a^2}{a}=a$, which is what we expected. On the other hand, options (C) and (D) equate to 0, which is clearly false.

 


	\item (\textit{AHSME 1965}) Point $E$ is selected on side $AB$ of $\triangle{ABC}$ in such a way that $AE: EB = 1: 3$ and point $D$ is selected on side $BC$ 
such that $CD: DB = 1: 2$. The point of intersection of $AD$ and $CE$ is $F$. Then $\frac {EF}{FC} + \frac {AF}{FD}$ is: 


 $\textbf{(A)}\ \frac {4}{5} \qquad  \textbf{(B) }\ \frac {5}{4} \qquad  \textbf{(C) }\ \frac {3}{2} \qquad  \textbf{(D) }\ 2\qquad \textbf{(E) }\ \frac{5}{2}$


 
 \textbf{Answer: }C



\textbf{Solution}


                    \begin{figure}[H]
                        centering
                        \includegraphics[width=0.5\textwidth]{Solution/AHSME_I_1965_37_Solution_Figure_1.png}
                    \end{figure}
                    

 We use \href{https://artofproblemsolving.com/wiki/index.php?title=Mass\_points}{mass points} for this problem. Let $\text{m} A$ denote the mass of point $A$.Rewrite the expression we are finding as 
 \[\frac{EF}{FC} + \frac{AF}{FD} = \frac{FE}{FC} + \frac{FA}{FD} = \frac{\text{m} C}{\text{m} E} + \frac{\text{m} D}{\text{m} A}\]Now, let $\text{m} C = 2$. We then have $2 \cdot 1 = \text{m} B \cdot 2$, so $\text{m} B = 1$, and $\text{m} D = 2+1 = 3$We can let $\text{m} A = 3$. We have $\text{m} E = \text{m} A + \text{m} B = 3+1 = 4$From here, substitute the respective values to get

 
 \[\frac{\text{m} C}{\text{m} E} + \frac{\text{m} D}{\text{m} A} = \frac{2}{4} + \frac{3}{3} = \frac{1}{2} + 1 = \frac{3}{2}\]This answer corresponds to $\fbox{\textbf{(C)}}$.

 


	\item (\textit{AHSME 1965}) $A$ takes $m$ times as long to do a piece of work as $B$ and $C$ together; $B$ takes $n$ times as long as $C$ and $A$ together; 
and $C$ takes $x$ times as long as $A$ and $B$ together. Then $x$, in terms of $m$ and $n$, is: 


 $\textbf{(A)}\ \frac {2mn}{m + n} \qquad  \textbf{(B) }\ \frac {1}{2(m + n)} \qquad  \textbf{(C) }\ \frac{1}{m+n-mn}\qquad \textbf{(D) }\ \frac{1-mn}{m+n+2mn}\qquad \textbf{(E) }\ \frac{m+n+2}{mn-1}$


 
 \textbf{Answer: }E



\textbf{Solution 1}

Let $a$, $b$, and $c$ be the speeds at which $A$, $B$ and $C$ work, respectively. Also, let the piece of work be worth one unit of work. Then, using the information from the problem along with basic \href{https://artofproblemsolving.com/wiki/index.php?title=Rate}{rate} formulas, we obtain the following equations:\begin{align*}\frac{1}{a}&=m*\frac{1}{b+c} \\\frac{1}{b}&=n*\frac{1}{a+c} \\\frac{1}{c}&=x*\frac{1}{a+b}\end{align*}
These equations can be rearranged into the following:\begin{align*}\text{(i) } ma&=b+c \\\text{(ii) } nb&=a+c \\\text{(iii) } xc&=a+b \\\end{align*}
Solving for $a$ in equation (i) gives us $a=\frac{b+c}{m}$. Substituting this expression for $a$ into equation (ii) yields:\begin{align*}nb&=\frac{b+c}{m}+c \\mnb&=b+c+mc \\(mn-1)b&=(m+1)c \\b&=\frac{(m+1)c}{mn-1} \end{align*}
Finally, substituting our expressions for $a$ and $b$ into equation (iii) yields our final answer:\begin{align*}xc&=\frac{b+c}{m}+\frac{(m+1)c}{mn-1} \\&=\frac{\frac{(m+1)c}{mn-1}+c}{m}+\frac{(m+1)c}{mn-1} \\&=\frac{c}{m}(\frac{m+1+mn-1}{mn-1}+\frac{m^2+m}{mn-1}) \\&=\frac{c}{m}(\frac{m^2+mn+2m}{mn-1}) \\&=c(\frac{m+n+2}{mn-1})\end{align*}
Thus, $x=\boxed{\textbf{(E) }\frac{m+n+2}{mn-1}}$.

 


 



\textbf{Solution 2 (Answer choices)}

If we let $A$, $B$, and $C$ work at the same speed, then it is clear that $m=n=x=2$. After plugging in $m=n=2$ into all of the answer choices, we see that the only two choices which give a value $x=2$ are choices (A) and (E). Now suppose that $A$ and $B$ work at the same speed, but $C$ does no work at all. Then, $m=n=1$, but $x$ is undefined. After plugging in $m=n=1$ into choices (A) and (E), we see that the only choice which is undefined is choice $\boxed{\textbf{(E) }\frac{m+n+2}{mn-1}}$.  

 


 


	\item (\textit{AHSME 1965}) A foreman noticed an inspector checking a $3$"-hole with a $2$"-plug and a $1$"-plug and suggested that two more gauges 
be inserted to be sure that the fit was snug. If the new gauges are alike, then the diameter, $d$, of each, to the nearest hundredth of an inch, is: 


 $\textbf{(A)}\ .87 \qquad  \textbf{(B) }\ .86 \qquad  \textbf{(C) }\ .83 \qquad  \textbf{(D) }\ .75 \qquad  \textbf{(E) }\ .71$


 
 \textbf{Answer: }B



\textbf{Solution 1}


                    \begin{figure}[H]
                        centering
                        \includegraphics[width=0.5\textwidth]{Solution/AHSME_I_1965_39_Solution_Figure_1.png}
                    \end{figure}
                    

 Let the center of the $3$" circle be $O$, that of the $2$" circle be $A$, that of the $1$" circle be $B$, and those of the circles of unknown radius (let their radii have length $r$) be $C$ and $D$, as in the diagram. Also, in this problem, no two circles share a center, so let the circles be named by their corresponding centers (so, the $3$" circle is circle $O$, etc.). Extend $\overline{OC}$ past $C$ to intersect circle $O$ at point $E$. Because circle $O$ has radius $\frac{3}{2}$ and circle $A$ has radius $1$, $OA=\frac{3}{2}-1=\frac{1}{2}$. Likewise, because $B$ has radius $\frac{1}{2}$, $OB=\frac{3}{2}-\frac{1}{2}=1$. Thus, $AB=\frac{3}{2}$. Furthermore, because the line connecting the centers of two tangent circles goes through their point of tangency, $AC=r+1$ and $BC=r+\frac{1}{2}$. Because $OE=\frac{3}{2}$ and $CE=r$, $OC=\frac{3}{2}-r$. With this information, we can now apply \href{https://artofproblemsolving.com/wiki/index.php?title=Stewart\%27s\_Theorem}{Stewart's Theorem} to $\triangle ABC$ to solve for $r$:\begin{align*}OB*OA*AB+OC^2*AB&=AC^2*OB+BC^2*OA \\4[(1)(\frac{1}{2})(\frac{3}{2})+(\frac{3}{2}-r)^2(\frac{3}{2})]&=4[(r+1)^2(1)+(r+\frac{1}{2})^2(\frac{1}{2})] \\3+6(\frac{9}{4}-3r+r^2)&=4(r^2+2r+1)+2(r^2+r+\frac{1}{4}) \\3+\frac{27}{2}-18r+6r^2&=4r^2+8r+4+2r^2+2r+\frac{1}{2} \\3+\frac{27-1}{2}-4+6r^2-6r^2&=18r+8r+2r \\3+13-4&=28r \\28r&=12 \\r&=\frac{3}{7}\end{align*}
Because the question asks for the diameter of the circle, we calculate $2r=\frac{6}{7}\approx \boxed{\textbf{(B) }0.86}$.

 



\textbf{Solution 2 (Overkill?)}

This is Descartes' formula for four mutually tangent circles with radii $r_1,r_2,r_3,r_4$

 Let $k_i$ denote $\frac{1}{r_i}$. This is the curvature.

 Note: A negative radius means that the other circles are internally tangent to it.

 Then,

 $(k_1+k_2+k_3+k_4)^2=2(k_1^2+k_2^2+k_3^2+k_4^2)$

 First, scale the problem up by two, so the radii are 1",2",3", and d"

 Next, let $k$ be $frac{1}{d}$

 Using the formula, we get:

 $\left(\frac{7}{6}+k\right)=2(\frac{49}{36}+k^2)$

 $\frac{49}{36}+\frac{7}{3}k+k^2=\frac{98}{36}+2k^2$

 (You'll see why I didn't simplify $\frac{98}{36}$)

 $k^2-\frac{7}{3}-\frac{49}{36}=0$

 Using the quadratic formula:

 $k=\frac{\frac{7}{3}\pm\sqrt{\frac{49}{9}-4\cdot1\cdot\frac{49}{36}}}{2}$

 $k=\frac{\frac{7}{3}\pm\sqrt{\frac{49}{9}-\frac{49}{9}}}{2}$

 $k=\frac{\frac{7}{3}\pm\sqrt{0}}{2}$

 $k=\frac{\frac{7}{3}\pm0}{2}$

 $k=\frac{\frac{7}{3}}{2}$

 $k=\frac{7}{6}$

 So the diameter $d=\frac{6}{7}$, which is about $\boxed{(\textbf{B}).86}$

 


	\item (\textit{AHSME 1965}) Let $n$ be the number of integer values of $x$ such that $P = x^4 + 6x^3 + 11x^2 + 3x + 31$ is the square of an integer. Then $n$ is: 


 $\textbf{(A)}\ 4 \qquad  \textbf{(B) }\ 3 \qquad  \textbf{(C) }\ 2 \qquad  \textbf{(D) }\ 1 \qquad  \textbf{(E) }\ 0$


 
 \textbf{Answer: }D



\textbf{Solution}

First, we wish to factor $P$ into a more manageable form. From the beginning of $P$, we notice $x^4+6x^3$, which gives us the idea to use $(x^2+3x)^2=x^4+6x^3+9x^2$. 

 This gives us 
 \[P=(x^2+3x)^2+2x^2+3x+31\]This is not useful, but it gives us a place to start from. 

 We can then try $(x^2+3x+1)=x^4+6x^3+11x^2+6x+1$. 
 \[P=(x^2+3x+1)^2-3(x-10)\]This is much more useful, as it moves all non-linear terms inside of a squared expression.

 We can then say $P=(x^2+3x+1)^2-3(x-10)=a^2$, where $a^2$ is the square of an integer mentioned on the problem. Right from here, we can set $x=10$, which cancels out the $3(x-10)$, giving $(100+30+1)^2=a^2$. This gives us one solution, $x=10$, $a=131$.

 
We can then rearrange the expression, giving us $(x^2+3x+1)^2-a^2=3(x-10)$. Factoring using difference of squares, we obtain 
 \[(x^2+3x+1+|a|)(x^2+3x+1-|a|)=3(x-10)\]

 We can then state that when $x$ is greater than $3$ and less than $-9$, $x^2+3x+1$ will be greater than $|3x-10|$. This is obtained by setting $x^2+3x+1>|3x-10|$ and then solving the inequality. We can then conclude that $x^2+3x+1+|a|>|3x-10|$. 

 Next, we claim that $x^2+3x+1-|a| \geq 1$ or $x^2+3x+1-|a| \leq -1$ when $x \neq 10$. We can prove this by first noting that since $x$ and $a$ are integers, $x^2+3x+1-|a|$ is an integer. Next, we shall assume that $x^2+3x+1-|a|=0$. Solving this and plugging back into the original equation, we obtain $(2|a|)(0)=3(x-10)$. Solving we obtain $x=10$, which is a contraction to $x \neq 10$. Therefore, $x^2+3x+1-|a| \neq 0$ and $x^2+3x+1-|a| \geq 1$ or $x^2+3x+1-|a| \leq -1$.

 Finally, we can go back to the equation 
 \[(x^2+3x+1+|a|)(x^2+3x+1-|a|)=3(x-10)\]We note that since $(x^2+3x+1+|a|)$ is larger than $3(x-10)$, in order for there to be solutions, $(x^2+3x+1-|a|)$ must be in the range $(-1,1)$. However, this contradicts what was proven earlier, so when $x \neq 10$ and $x < -9$ or $x > 3$, there are no solutions for $x$. 

 
Now, all that remains to be checked are values of $x$ between $-9$ and $3$. Using brute force and checking each value individually, we can assert that there are no such solutions for $x$, leaving us with only $1$ solution, $x=10$. Therefore, the answer is $\boxed{(D) }$

 Solution by treetor10145

 


\end{enumerate}